\section{Expansion}
\label{sec:expansion}

Exponential shell growth transforms both cellular-automaton reversibility and wave propagation on $\{3,4,3,4\}$. Because over $94.5\%$ of every large ball lies on its single outermost shell, the Euclidean Moore--Myhill equivalence fails, post-hoc injectivity testing for arbitrary local rules becomes unsolvable, and every random walk evacuates outward at a Lean~4-certified exponential rate exceeding $13/25$ per step.

\subsection{Non-Amenability}
\label{subsec:non-amenability}

On a flat lattice $\mathbb{Z}^d$, the boundary of a ball of radius $r$ grows as $O(r^{d-1})$ against volume $\Theta(r^d)$, so the surface-to-volume ratio $|\partial B_r|/|B_r|$ vanishes as $r \to \infty$. That amenability property drives the Moore--Myhill Garden-of-Eden theorem~\cite{moore1962,myhill1963,bartholdi2010}: a cellular automaton on an amenable group is globally surjective iff its restriction to finite configurations is injective. On hyperbolic tilings, boundary area stays a strictly positive fraction of volume. On the planar heptagrid $\{7,3\}$, Margenstern showed that $|\partial B_r|/|B_r| > 1/2$ for all radii, breaking Moore--Myhill so that injective-but-non-surjective and surjective-but-non-injective automata coexist~\cite{margenstern2007hedlund,margenstern2008injectivity}. Extending Kari's planar tiling reduction~\cite{kari1990,kari1994} via an almost plane-filling path across interwoven triangles, Margenstern further proved that deciding global injectivity of an arbitrary cellular automaton on $\{7,3\}$ is undecidable~\cite{margenstern2008injectivity}.

On $\{3,4,3,4\}$, boundary dominance is far sharper: by Theorem~\ref{thm:cubic-growth} and~\eqref{eq:surface-vol}, $|\partial B_r|/|B_r| \to 0.945292$. Any finite patch of $n$ cells touches an outer perimeter nearly as large as its entire interior. Preimage collisions in an unstructured local rule table can therefore hide at arbitrarily large radii, making empirical injectivity checks futile. At the same time, strong non-amenability forces a wide spectral gap on the $24$-regular adjacency operator, ensuring that diffuse excitations decay exponentially fast rather than stalling near an active region.

\subsection{Spectral Bounds}
\label{subsec:spectral-bounds}

A normalized random-walk spectral gap above $13/25$ means that every vector in $\ell^2(W(\Gamma_{24}))$ shrinks in norm by more than a factor of $25/12$ on every step. Let $A$ be the $24$-regular adjacency operator on $\ell^2(W(\Gamma_{24}))$, with spectral radius $\rho = \|A\|_{2 \to 2}$ and normalized random-walk gap
\begin{equation}
\label{eq:gap-def}
\lambda_{\mathrm{gap}} = 1 - \frac{\rho}{24}.
\end{equation}
By Kesten's criterion~\cite{kesten1959} (a finitely generated group is non-amenable iff its simple random walk has spectral radius strictly below $1$), non-amenability implies $\rho < 24$. Using reciprocal edge weights on Cartier--Foata suffix automata and Markov Rayleigh--Ritz test vectors, we collapse $\rho$ from a Depth-1 bracket of width below $0.8764$ to a Lean~4-certified Depth-2 bracket of width below $0.086$ (over $10.2\times$ tighter, certifying $\lambda_{\mathrm{gap}} > 13/25$) and an exact big-integer Depth-5 bracket of width below $0.0047$ (over $187\times$ tighter)~\cite{pollard2026repo}.

\paragraph{Upper bounds via reciprocal edge weights.}
Following the Gabber--Galil--Nagnibeda Schur test~\cite{nagnibeda1997}, assign a positive weight $w(u \to v) > 0$ to each directed edge of a symmetric finite graph $(V, A)$ subject to the reciprocal condition
\begin{equation}
\label{eq:reciprocal-schur}
w(u \to v)\,w(v \to u) = 1 \qquad \text{for every edge } u \sim v.
\end{equation}
Because the arithmetic--geometric mean inequality gives
\begin{equation}
\label{eq:ag-mean}
2\,|x_u x_v| \le w(u \to v)\,x_u^2 + w(v \to u)\,x_v^2,
\end{equation}
summing over all edges bounds the operator norm without requiring any curl-free vertex potential:
\begin{equation}
\label{eq:schur-quad}
|\langle x, A x \rangle| \le \Bigl(\max_{u \in V} \sum_{v \sim u} w(u \to v)\Bigr)\,\langle x, x \rangle.
\end{equation}
At Depth~1 on the four Cartier--Foata strata $(C_0, C_1, C_2, C_3)$, assigning three positive rationals $(p_1, p_2, p_3)$ to upward edges entering stratum $C_k$ and their exact reciprocals $(p_1^{-1}, p_2^{-1}, p_3^{-1})$ to reverse downward edges collapses the row sums in~\eqref{eq:schur-quad} to four closed forms:
\begin{align}
S_0 &= 24\,p_1, & S_1 &= p_1^{-1} + 15\,p_1 + 8\,p_2, \label{eq:schur-s01}\\
S_2 &= 2\,p_2^{-1} + 11\,p_1 + 8\,p_2 + 3\,p_3, & S_3 &= 3\,p_3^{-1} + 9\,p_1 + 6\,p_2 + 6\,p_3. \label{eq:schur-s23}
\end{align}
To sharpen past the Depth-1 ceiling, refine each vertex state from its single rightmost clique $D_R(w)$ to the $W(F_4)$-orbit of its last $d_{\mathrm{suf}}$ Cartier--Foata commuting blocks $\mathrm{St}_{d_{\mathrm{suf}}}(w)$. Tracking these multi-block suffixes expands the finite quotient automaton from $4$ states at $d_{\mathrm{suf}} = 1$ to $19$ states at $d_{\mathrm{suf}} = 2$ and $34{,}523$ states at $d_{\mathrm{suf}} = 5$ (through $169$ and $2{,}197$ at depths $3$ and $4$). Labeling every upward edge $u \to v$ by the $W(F_4)$-orbit upward-end flag $[\mathrm{St}_{d_{\mathrm{suf}}}(v), b]_{W(F_4)}$ of the entering generator $b \in D_R(v)$ ($23$ flags at $d_{\mathrm{suf}} = 2$, rising to $43{,}312$ at $d_{\mathrm{suf}} = 5$) fixes both directions of every edge from the child suffix alone, enforcing exact reciprocity automatically across all finite subgraphs while equalizing every row sum using quantized rational weights with denominator scale $q_{\mathrm{den}}$.

\paragraph{Lower bounds via Cartier--Foata Markov Rayleigh--Ritz.}
At Depth~1, finite-support stratum vectors
\begin{equation}
\label{eq:rr-ansatz}
\phi(r, k) = a_k\,r\,(R + 1 - r)\,11^{r-1}\,47^{R-r}
\end{equation}
use the rational scaling $(11/47)^r \approx \lambda_{3434}^{-r/2}$ to balance exponential shell growth and compress $\langle x, A x \rangle$ identically onto the $3R$-dimensional stratum space, yielding $\rho > 11{,}049/1{,}000$. Because spherical stratum weights average over local two-clique correlations, we lift the Rayleigh ansatz to a $W(F_4)$-invariant Cartier--Foata Markov product measure that weights each adjacent clique transition $(C_t, C_{t+1})$ by its $57$-orbit inner-product class (and at $2$-step depth by $4{,}440$ clique triples). Right-multiplying a Cartier--Foata word by a facet generator $b$ slides $b$ leftward only through generators commuting with $b$ inside the $3$-cube $\mathrm{Adj}(b)$ ($20$ nonempty cliques), reducing the Rayleigh numerator over Cartier--Foata block heights $h \le H_{\max}$ with zero dropped negative terms (safely truncating slide depth at $M = 11$ in Lean~4) and certifying exact integer-vector lower bounds.

\begin{theorem}[Certified Depth-2 Spectral Bracket and Depth-5 Collapse]
\label{thm:spectral-bracket}
Let $A$ be the adjacency operator on a finite bipartite subgraph matching the $W(F_4)$-orbit Cartier--Foata suffix counts and upward-end flag multiplicities of $\{3,4,3,4\}$, and let $\rho = \|A\|_{2\to 2}$ be the infinite-graph spectral radius obtained by density of finitely supported vectors in $\ell^2(W(\Gamma_{24}))$.
\begin{enumerate}
\item \textbf{Depth-1 baseline (Lean~4):} Substituting $(p_1, p_2, p_3) = (1/7,\, 1/3,\, 4/7)$ into~\eqref{eq:schur-s01}--\eqref{eq:schur-s23} gives $\rho \le 335/28 < 12$ and $\lambda_{\mathrm{gap}} \ge 337/672 > 1/2$, refined at Depth~1 (width $< 0.8764$) to
\begin{equation}
\label{eq:rho-depth1}
\frac{11{,}049}{1{,}000} < \rho\bigl(\{3,4,3,4\}\bigr) \le \frac{3{,}325{,}977{,}059}{278{,}900{,}000}.
\end{equation}
\item \textbf{Depth-2 kernel collapse (Lean~4):} Over the $19$-state Depth-2 suffix automaton ($23$ reciprocal flags, denominator scale $q_{\mathrm{den}} = 10^4$) and the $57$-orbit $1$-step Cartier--Foata Markov witness over block heights $h \le H_{\max} = 120$, Lean~4 certifies
\begin{equation}
\label{eq:rho-bracket}
\frac{913}{80} < \rho\bigl(\{3,4,3,4\}\bigr) \le \frac{387{,}141{,}227}{33{,}670{,}000} < \frac{115}{10},
\qquad
\lambda_{\mathrm{gap}} \ge \frac{420{,}938{,}773}{808{,}080{,}000} > \frac{25}{48} > \frac{13}{25} = 0.52,
\end{equation}
shrinking the kernel-certified interval width to $57{,}647{,}040/673{,}400{,}000 < 0.086$ (more than $10.2\times$ tighter than Depth~1).
\item \textbf{Depth-5 exact big-integer collapse (companion):} At Depth~5 ($34{,}523$ states, $43{,}312$ reciprocal flags at $q_{\mathrm{den}} = 10^6$) paired with a concrete finite-support $2$-step ($4{,}440$-triple) Cartier--Foata Markov witness over block heights $h \in \{10,\dots,420\}$, exact big-integer arithmetic certifies
\begin{equation}
\label{eq:rho-depth5}
11.439047 \le \rho\bigl(\{3,4,3,4\}\bigr) \le 11.443713, \qquad \lambda_{\mathrm{gap}} \ge 0.5231,
\end{equation}
narrowing the certified interval width to $0.004666$ (over $187\times$ tighter than Depth~1).
\end{enumerate}
\end{theorem}

Table~\ref{tab:spectral-summary} and Figure~\ref{fig:spectral-growth} display all three nested certificates beside the multi-dimensional capacity curve. In Lean~4, the $19$-state Depth-2 automaton, all $23$-flag reciprocities, the finite-subgraph Schur theorem, the $57$-orbit $Q_3$ cascade reduction, and every rational inequality in~\eqref{eq:rho-bracket} are verified unconditionally~\cite{pollard2026repo}. Passing from finite subgraphs and Cartier--Foata words to the infinite operator norm on $\ell^2(W(\Gamma_{24}))$ uses Tits's Cartier--Foata normal-form bijection on right-angled Coxeter groups together with standard density of finitely supported vectors.

\begin{table}[tb]
\centering
\caption{Exact rational and big-integer spectral certificates for the $24$-regular dual graph of $\{3,4,3,4\}$ across Cartier--Foata depths $1$, $2$, and $5$~\cite{pollard2026repo}. Decimal bounds round strictly outward.}
\label{tab:spectral-summary}
\footnotesize
\setlength{\tabcolsep}{4.5pt}
\begin{tabular}{@{}lllll@{}}
\toprule
Level & Automaton / witness & Exact bound & Safe decimal & Verification \\
\midrule
Depth~1 lower & $R \ge 36$, $(17,26,45)$ & $> 11{,}049/1{,}000$ & $> 11.0490$ & Lean~4 \\
Depth~1 upper & $4$ states, $3$ flags & $\le 3{,}325{,}977{,}059/278{,}900{,}000$ & $\le 11.9254$ & Lean~4 \\
Depth~2 lower & $57$-orbit CF ($H_{\max}=120$) & $> 913/80$ & $> 11.4125$ & Lean~4 \\
Depth~2 upper & $19$ states, $23$ flags ($q_{\mathrm{den}}=10^4$) & $\le 387{,}141{,}227/33{,}670{,}000$ & $< 11.4982$ & Lean~4 \\
Depth~2 gap & $\lambda_{\mathrm{gap}} = 1 - \rho/24$ & $\ge 420{,}938{,}773/808{,}080{,}000$ & $> 0.5209$ & Lean~4 \\
Depth~5 lower & $4{,}440$-triple CF ($h \le 420$) & Exact $\mathbb{Z}$ Rayleigh ratio & $\ge 11.439047$ & BigInt comp. \\
Depth~5 upper & $34{,}523$ states, $43{,}312$ flags & Exact $q_{\mathrm{den}}=10^6$ Schur & $\le 11.443713$ & BigInt comp. \\
\bottomrule
\end{tabular}
\end{table}

\begin{figure}[tb]
\centering
\begin{tikzpicture}[scale=0.95, every node/.style={font=\small}]
  % Subtle center divider only (pure white background, zero gray fill)
  \draw[twZincRule, thin] (7.68,-0.62) -- (7.68,4.42);

  % Left panel: log10 cumulative capacity vs depth d=0..5 (x = 0.0 .. 7.45)
  \node[anchor=west, font=\small\bfseries, text=twZincInk] at (0.0,4.26) {Cumulative capacity $\log_{10} V(d)$};

  % Axes: x in [0.80, 4.90] for d=0..5 (0.82 cm per step), y in [0.45, 3.45] for log10 in [0, 7] (0.42 cm per unit)
  \draw[-{Stealth[length=4pt]}, twZincInk, thin] (0.80,0.45) -- (5.16,0.45);
  \draw[-{Stealth[length=4pt]}, twZincInk, thin] (0.80,0.45) -- (0.80,3.45);
  \node[anchor=south, font=\scriptsize, text=twZincInk] at (0.80,3.48) {$\log_{10} V$};
  \node[font=\scriptsize, text=twZincInk] at (2.90,-0.44) {graph depth $d$};

  % Ticks x: d = 0..5
  \foreach \d in {0,1,2,3,4,5} {
    \pgfmathsetmacro{\xx}{0.80 + \d*0.82}
    \draw[twZinc, thin] (\xx,0.38) -- (\xx,0.52);
    \node[below, font=\scriptsize, text=twZinc] at (\xx,0.34) {$\d$};
  }
  % Ticks y: 0, 2, 4, 6
  \foreach \y in {0,2,4,6} {
    \pgfmathsetmacro{\yy}{0.45 + \y*0.42}
    \draw[twZinc, thin] (0.73,\yy) -- (0.87,\yy);
    \node[left, font=\scriptsize, text=twZinc] at (0.69,\yy) {$\y$};
  }

  % Curve 1: 4D {3,4,3,4} cusp ball V_d = 1, 19, 349, 6383, 116681, 2132795
  \draw[twEmerald, very thick] (0.80,0.450) -- (1.62,0.987) -- (2.44,1.518) -- (3.26,2.048) -- (4.08,2.578) -- (4.90,3.108);
  \foreach \x/\y in {0.80/0.450, 1.62/0.987, 2.44/1.518, 3.26/2.048, 4.08/2.578, 4.90/3.108} {
    \fill[twEmerald] (\x,\y) circle (2.2pt);
  }
  \node[anchor=west, font=\scriptsize\bfseries, text=twEmerald] at (5.00,3.11) {$\{3,4,3,4\}$ cusp};

  % Curve 2: 3D {5,3,4} exact RACG ball V_d = 1, 13, 115, 927, 7329, 57741
  \draw[twBlue, thick] (0.80,0.450) -- (1.62,0.918) -- (2.44,1.315) -- (3.26,1.696) -- (4.08,2.073) -- (4.90,2.450);
  \foreach \x/\y in {0.80/0.450, 1.62/0.918, 2.44/1.315, 3.26/1.696, 4.08/2.073, 4.90/2.450} {
    \fill[twBlue] (\x,\y) circle (2.0pt);
  }
  \node[anchor=west, font=\scriptsize, text=twBlue] at (5.05,2.45) {$\{5,3,4\}$ ball};

  % Curve 3: Flat Z^4 l1 ball: 1, 9, 41, 129, 321, 681
  \draw[twViolet, thick, dashed] (0.80,0.450) -- (1.62,0.851) -- (2.44,1.127) -- (3.26,1.336) -- (4.08,1.503) -- (4.90,1.640);
  \foreach \x/\y in {0.80/0.450, 1.62/0.851, 2.44/1.127, 3.26/1.336, 4.08/1.503, 4.90/1.640} {
    \fill[twViolet] (\x,\y) circle (2.0pt);
  }
  \draw[twViolet, thin] (4.94,1.66) -- (5.08,1.88);
  \node[anchor=west, font=\scriptsize, text=twViolet] at (5.08,1.89) {flat $\mathbb{Z}^4$ ball};

  % Curve 4: 2D {7,3} heptagrid ball: 1, 8, 29, 85, 232, 617
  \draw[twZinc, thick, dotted] (0.80,0.450) -- (1.62,0.829) -- (2.44,1.064) -- (3.26,1.260) -- (4.08,1.443) -- (4.90,1.622);
  \draw[twZinc, thin] (4.94,1.60) -- (5.08,1.32);
  \node[anchor=west, font=\scriptsize, text=twZinc] at (5.08,1.30) {2D $\{7,3\}$ ball};

  % Right panel: Nested Spectral Collapse on [0, 24] and Zoom on [11.0, 12.0] (x = 7.95 .. 15.25)
  \node[anchor=west, font=\small\bfseries, text=twZincInk] at (8.05,4.26) {Nested $\{3,4,3,4\}$ spectral collapse};

  % Top number line for rho in [0, 24] at y = 2.78 (x = 8.25 to 14.95, 6.70 cm)
  \draw[twZincInk, thick] (8.25,2.78) -- (14.95,2.78);
  \foreach \v in {0,6,12,18,24} {
    \pgfmathsetmacro{\xv}{8.25 + \v*6.70/24}
    \draw[twZincInk, thin] (\xv,2.68) -- (\xv,2.88);
    \node[below, font=\scriptsize, text=twZinc] at (\xv,2.62) {$\v$};
  }

  % Bracket highlight on top bar: D1 11.049..11.9254 -> 11.334..11.579; D2 11.4125..11.4982 -> 11.436..11.460
  \draw[twZincRule, line width=3.2pt] (11.334,2.78) -- (11.579,2.78);
  \draw[twViolet, line width=4.2pt] (11.432,2.78) -- (11.465,2.78);

  % Midpoint 12 dashed vertical line and left-aligned label cleanly separated from gap label
  \draw[twPink, dashed, thin] (11.60,2.68) -- (11.60,3.70);
  \node[anchor=east, font=\scriptsize, text=twPink] at (11.48,3.60) {$\rho = 12$};

  % Gap span arrow cleanly above top bar (y = 3.18), label right of pink line
  \draw[{Stealth[length=4pt]}-{Stealth[length=4pt]}, twEmerald, thick] (11.66,3.18) -- (14.95,3.18);
  \node[above, font=\scriptsize\bfseries, text=twEmerald] at (13.35,3.24) {$\lambda_{\mathrm{gap}} > 0.5209$};

  % Zoom callout lines terminating strictly ABOVE the Depth-1 text label (y=2.24 -> y=1.88, ZERO text crossings)
  \draw[twZincRule, thin, dashed] (11.32,2.24) -- (8.45,1.88);
  \draw[twZincRule, thin, dashed] (11.60,2.24) -- (14.75,1.88);

  % Zoomed inset below on rho in [11.0, 12.0] (x = 8.45 to 14.75, 6.30 cm per unit)
  % Axis at y = 0.34
  \draw[twZincInk, thick] (8.45,0.34) -- (14.75,0.34);
  \foreach \v/\lbl in {0/11.0, 0.25/11.25, 0.50/11.50, 0.75/11.75, 1.00/12.0} {
    \pgfmathsetmacro{\xz}{8.45 + \v*6.30}
    \draw[twZincInk, thin] (\xz,0.26) -- (\xz,0.42);
    \node[below, font=\scriptsize, text=twZinc] at (\xz,0.20) {$\lbl$};
  }

  % Row 1 (y=1.42): Depth-1 [11.049, 11.9254] -> x = 8.759 .. 14.280
  \draw[twBlue, line width=2.8pt] (8.76,1.42) -- (14.28,1.42);
  \fill[twBlue] (8.76,1.42) circle (2.0pt);
  \fill[twBlue] (14.28,1.42) circle (2.0pt);
  \node[above, font=\scriptsize, text=twBlue, fill=white, inner sep=1.2pt] at (11.52,1.48) {Depth~1 (Lean~4): $(11.049,\, 11.9254)$, width $< 0.8764$};

  % Row 2 (y=1.00): Depth-2 (11.4125, 11.4982] -> x = 11.049 .. 11.589
  \draw[twViolet, line width=3.4pt] (11.05,1.00) -- (11.59,1.00);
  \fill[twViolet] (11.05,1.00) circle (2.2pt);
  \fill[twViolet] (11.59,1.00) circle (2.2pt);
  \node[anchor=west, font=\scriptsize\bfseries, text=twViolet] at (11.72,1.00) {Depth~2 (Lean~4): width $< 0.086$};

  % Row 3 (y=0.62): Depth-5 [11.439047, 11.443713] (outward [11.4390, 11.4438]) -> x = 11.216 .. 11.245
  \draw[twEmerald, line width=4.0pt] (11.21,0.62) -- (11.26,0.62);
  \node[anchor=west, font=\scriptsize\bfseries, text=twEmerald] at (11.40,0.62) {Depth~5 (BigInt): $[11.4390,\, 11.4438]$};

  \node[font=\scriptsize, text=twZinc] at (11.60,-0.44) {Zoomed interval $\rho \in [11.0,\, 12.0]$ ($> 187\times$ collapse)};
\end{tikzpicture}
\caption{Quantitative expansion on $\{3,4,3,4\}$. Left: exact logarithmic cumulative capacity $\log_{10} V(d)$ through depth $d = 5$, comparing $\{3,4,3,4\}$ cusp balls against exact 3D $\{5,3,4\}$ RACG balls, 2D $\{7,3\}$ balls, and flat $\mathbb{Z}^4$ $\ell_1$-balls. Right: nested spectral collapse from Depth~1 (width $< 0.8764$) to Lean~4-certified Depth~2 ($(913/80,\, 387{,}141{,}227/33{,}670{,}000]$, width $< 0.086$, $\lambda_{\mathrm{gap}} > 13/25$) and exact big-integer Depth~5 ($[11.439047,\, 11.443713]$, width $< 0.0047$).}
\label{fig:spectral-growth}
\end{figure}

From a cellular-automata standpoint, Theorem~\ref{thm:spectral-bracket} underscores a clear structural requirement. With $\lambda_{\mathrm{gap}} > 13/25$ and $|\partial B_r|/|B_r| \to 0.9453$, unstructured local rules sit deep inside the non-amenable regime where global injectivity and surjectivity decouple and post-hoc injectivity verification is undecidable. Constructing provably reversible universal computation on $\{3,4,3,4\}$ therefore calls for a local algebraic factorization that enforces bijectivity by construction.
